\begin{abstract}
A common machine-learning technique adds an \emph{auxiliary loss}
that pushes a network's internal features toward a known target
direction in feature space---for example, the difference between
positive and negative sentiment, or between calming and disturbing
emotion. We test $25$ such recipes (knowledge distillation,
representational-similarity analysis, contrastive losses, low-rank
adapters, channel-targeted brain-region losses) on
FACED~\citep{chen2023faced}, a public dataset of EEG recordings from
$123$ subjects watching $28$ emotional video clips, with $9$ emotion
classes. None help: $16$ of the $25$ produce statistically significant
accuracy decrements, and $0$ produce gains. We trace this to a
regularity we call \emph{saturation}: once a network has converged on
the emotion classification task from labels alone, its features
already encode the target concept; the auxiliary loss then adds
gradient noise rather than information. Saturation is distinct from
overfitting (memorising the training set); it is a property of how
task and auxiliary supervision interact once both target the same
direction in feature space. The same picture suggests a positive
prescription: ensemble across the seed-specific differences in the
features, rather than supervise the part of the features that has
already saturated. Following this prescription on FACED yields a new
state of the art at $\mathbf{0.6948}$ balanced accuracy---a $+10.5\%$
relative improvement over EMOD~\citep{chen2025emod}, the previous SOTA
on FACED, at $0.6287$. Surprisingly, the language-model-derived class
prototypes used by the knowledge-distillation step in our recipe are
interchangeable with random orthonormal vectors ($\Delta\le 0.003$
balanced accuracy): the SOTA gain comes from the \emph{shape} of the
loss (a 9-class structure imposed on the features), not from the
\emph{content} of the prototypes. To probe what the network has
saturated on---and to demonstrate it is a real semantic axis, not a
noise dimension---we extract one direction explicitly. We write nine
short emotion-evocative stories (one per FACED class), pass them
through a language model (Qwen3-4B), and take the first principal
component of the language model's late-layer features across the nine
stories. We call this direction the \emph{valence axis}. The same
direction predicts movie-review sentiment zero-shot on
SST-2~\citep{socher2013recursive} at AUC $0.832$, predicts the EEG
response averaged across $123$ FACED subjects at $r=0.87$, and emerges
spontaneously in $36$ EEG classifiers trained only on the $9$ emotion
labels. All claims replicate on a second public EEG dataset,
SEED-V~\citep{liu2022seedv} (Appendix~\ref{app:seedv-full}).
\end{abstract}
